\newpage
\subsection{Realisierung mit Bipolar- und FET-OPV}\label{sec:1-2}
\subsubsection{Aufgabenstellung}
Die Schaltung~\ref{fig:integrator-schaltung} soll sowohl mit dem Bipolar-OPV $A_3$ ($\mu$A741) als auch mit einem der FET-OPVs $A_1$/$A_2$ (TL081) aufgebaut werden, um deren jeweilige Biasströme zu bestimmen.

\subsubsection{Versuchsaufbau}
Die in Abbildung~\ref{fig:ibn} dargestellte Schaltung wird erweitert, sodass der Kondensator ($C$) über den Schalter ($S$) entladen werden kann. Die entsprechende Anpassung ist in Abbildung~2 dargestellt. Die Werte für ($R$) und ($C$) entsprechen weiterhin jenen aus Abschnitt~\ref{sec:1-1}.

Durch die virtuelle Masse am invertierenden Eingang $V_-$ kann der Widerstand aus Abb.~\ref{fig:integrator-schaltung} entfallen. Zur Bestimmung des Eingangsruhestroms genügt es, den negativen Biasstrom zu messen, da der nichtinvertierende Eingang des OPV auf Masse liegt.

\begin{figure}[h]
	\centering
	\begin{circuitikz}[scale=0.9, transform shape]
		% --- Op-amp A3 ---
		\node[op amp] (A3) at (0,0) {};
		\node at ($(A3.center)+(0.8,0.3)$) {\small $A_3$};

		% Supply arrows (+12V, -12V)
		\draw[->] ($(A3.center)+(0.2,0.4)$) -- +(0,0.5) node[above]{\small $+12$V};
		\draw[->] ($(A3.center)+(0.2,-0.4)$) -- +(0,-0.5) node[below]{\small $-12$V};

		% --- Output node ---
		\coordinate (out) at ($(A3.out)+(1.0,0)$);
		\draw (A3.out) -- (out);

		%capacitor and switch
		\coordinate (a-neg-mid) at ($(A3.-)+(-0.2,0)$);
		\coordinate (a-out-mid) at ($(A3.out)+(0.2,0)$);
		\coordinate (c-in) at ($(a-neg-mid)+(0,2.7)$);
		\coordinate (c-out) at ($(a-out-mid)+(0,3.2)$);
		\coordinate (switch-in) at ($(c-in)+(0,1.4)$);
		\coordinate (switch-out) at ($(c-out)+(0,1.4)$);
		\draw (c-in) to[C, l^={C},*-*] (c-out);
		\draw (switch-in) to[ospst, l^={$S_1$}] (switch-out);
		\draw (switch-in) -- (a-neg-mid);
		\draw (switch-out) -- (a-out-mid);
		\draw[->] ($(c-in)+(+0.5,-0.7)$) -- ($(c-out)+(-0.5,-0.7)$)
		node[midway,below] {$U_C$};

		% V_+
		\node[ground] at (A3.+) {};

		% --- Non-inverting input: voltage source G1 with U_os ---
		\coordinate (srcbottom) at (-2,-1.7);
		\node[ground] at ($(srcbottom)+(0,0.5)$) {};
		\draw (srcbottom) to[current source, i_={$I_{Bn}$}] ++(0,2.2) coordinate (srcnode);
		\draw (srcnode) -- (A3.-);

		% Output 
		\node[ocirc] at (out) {};
		\draw[->] ($(out.south)+(0,-0.2)$) -- ++(0,-0.7)
		node[midway,right] {$U_a$}
		coordinate (mid);
		\node[ocirc] (C2) at ($(mid)+(0,-0.2)$) {};
		\draw (C2.south) -- ++(0,-0.0) node[ground] {};

	\end{circuitikz}
	\caption{Messschaltung für $I_{Bn}$}\label{fig:ibn}
\end{figure}

\subsubsection{Messung}
Die Schaltung aus Abb.~\ref{fig:ibn} wird sowohl mit $A_1$ als auch mit $A_3$ aufgebaut. In Abb.~\ref{fig:ibn-a1-oszi} und Abb.~\ref{fig:ibn-a3-oszi} ist die jeweilige Ausgangsspannung dargestellt, während sich der Kondensator nach dem Öffnen des Schalters $S_1$ mit dem Strom $I_{Bn}$ auflädt.

\begin{figure}[h]
	\begin{minipage}{0.49\textwidth}
		\centering
		\includegraphics[width=1\textwidth]{oszi/12-IBn-A1.JPG}
		\caption{$U_{a}$ des FET-OPV $A_1$}\label{fig:ibn-a1-oszi}
	\end{minipage}\hfill
	\begin{minipage}{0.49\textwidth}
		\centering
		\includegraphics[width=1\textwidth]{oszi/12-IBn-A3.JPG}
		\caption{$U_{a}$ des Bipolar-OPV $A_3$}\label{fig:ibn-a3-oszi}
	\end{minipage}
\end{figure}

\subsubsection{Berechnung}
In der Schaltung~\ref{fig:integrator-schaltung} gilt $U_a = -U_C$ sowie $I_{Bn} = I_C$. Für den Kondensator gilt die bekannte Strom-Spannungsbeziehung $I_C = C\, \frac{dU_C}{dt}$. Damit ergibt sich

\begin{gather}
	I_{Bn} = C \frac{\Delta U_a}{\Delta t} \label{eq:ibn-oszi}
\end{gather}

Aus Abb.~\ref{fig:ibn-a1-oszi} und Abb.~\ref{fig:ibn-a3-oszi} sowie Gleichung~\ref{eq:ibn-oszi} ergeben sich folgende Werte:

\begin{table}[h]
	\centering
	\begin{tabular}{c|c|c|c}
		\textbf{OPV} & \textbf{Typ} & \textbf{$\frac{\Delta U_a}{\Delta t}$} & \textbf{$I_{Bn}$}         \\
		\hline
		$A_1$        & FET          & \SI{161.29}{\micro\volt\per\second}    & \SI{161.29}{\pico\ampere} \\
		$A_3$        & Bipolar      & \SI{40.32}{\milli\volt\per\second}     & \SI{40.32}{\nano\ampere}  \\
	\end{tabular}
	\caption{Messwerte der Operationsverstärker}
	\label{tab:op-messwerte}
\end{table}

\subsubsection{Erkenntnisse}
Vergleicht man die gemessenen Eingangsruheströme mit den typischen Datenblattangaben, so liegt der Wert von $I_{Bn}$ für den FET-OPV $A_1$ leicht über den üblichen \SI{30}{\pico\ampere}, während der gemessene Wert für den Bipolar-OPV $A_3$ etwas unterhalb der typischen \SI{80}{\nano\ampere} liegt. Beide Messwerte befinden sich jedoch in einem plausiblen Bereich.
